\section{State on $SE_2(3)$}

%%TODO: this sounds a bit clunky - invariant observer -> observable, a bit repetitive.
We define the state on the group manifold $SE_2(3)$, such that we can maintain invariant observer properties while keeping the system as observable as possible:

\begin{equation}
  \begin{gathered}
  X \triangleq
  \begin{bmatrix}
    \rotmat{W}{B} & \fv{W}{v}_{\mathcal{B}} & \fv{W}{p}_{\mathcal{B}} \\
    \mathbf{0}_{1\times 3} & 1 & 0 \\
    \mathbf{0}_{1\times 3} & 0 & 1
  \end{bmatrix}
  \in SE_2(3) \\
  \boldsymbol{\theta} = (\boldsymbol{b}_g, \boldsymbol{b}_a) \in \mathbb{R}^6
  \end{gathered}
  \label{eq:car_state}
\end{equation}

this has the inverse state also defined as:

\begin{equation}
  X^{-1} =
  \begin{bmatrix}
    R^\top & -R^\top \mathbf{v} & -R^\top \mathbf{p} \\
    \mathbf{0}_{1\times 3} & 1 & 0 \\
    \mathbf{0}_{1\times 3} & 0 & 1
  \end{bmatrix} \in SE_2(3)
  \label{eq:inverse_state}
\end{equation}

As the state is defined on a \textit{matrix lie group}, we also have to define the associated lie algebra state and its associated tangent vector. Both are defined below:

\begin{equation}
  \bsxi^\wedge = 
  \begin{bmatrix}
    \bsrho_{\bsphi}^\wedge & \bsrho_{\mathbf{v}} & \bsrho_{\mathbf{p}} \\
    0 & 0 & 0 \\
    0 & 0 & 0
  \end{bmatrix}
  \label{eq:lie_algebra}
\end{equation}

\begin{equation}
  \bsxi = 
  \begin{bmatrix} 
    \bsrho_{\bsphi} & \bsrho_{\mathbf{v}} & \bsrho_{\mathbf{p}} & \boldsymbol{\zeta}_g & \boldsymbol{\zeta}_a
  \end{bmatrix} \in \mathbb{R}^{15}
  \label{eq:tangent_vector}
\end{equation}

where the exponential on the state component of the tangent vector defined as
$\exp([\bsxi_{1:9}]^\wedge) = \hat{X}X^{-1}=\bseta_R$ (the right-invariant error), and $\boldsymbol{\zeta} = \hat{\boldsymbol{\theta}} - \boldsymbol{\theta}$.
